\subsection{Polytope representations and finite geometry}
\label{subsec:preliminaries-polytope-finite-geometry}

The thesis uses two equivalent finite descriptions of a full-dimensional
polytope \(K\subset\mathbb{R}^4\) with \(0\in\operatorname{int}(K)\). An
\(H\)-representation is an intersection of normalized halfspaces
\[
  K=\{x\in\mathbb{R}^4:\langle a_i,x\rangle\leq 1,\; i=1,\ldots,F\}.
\]
Here \(a_i=n_i/h_i\), where \(n_i\) is the outward unit normal and
\(h_i=h_K(n_i)>0\) is the corresponding height. The same data are the vertices
of the polar body
\[
  K^\circ=\operatorname{conv}\{a_1,\ldots,a_F\}
\]
when the inequalities are irredundant. Thus checking that a proposed dual-row
list really describes an \(F\)-facet polytope has two parts: \(0\) must lie in
the interior of \(\operatorname{conv}\{a_i\}\), which is equivalent to
boundedness of \(K\), and every \(a_i\) must be an extreme point of that convex
hull, which is equivalent to irredundancy of the normalized halfspace
description.

The finite vertex enumeration used in the implementation is the standard
four-dimensional halfspace algorithm. For every four-element set
\(I=\{i_1,i_2,i_3,i_4\}\), solve
\[
  \langle a_{i_j},x\rangle=1,
  \qquad j=1,\ldots,4.
\]
If the system is singular, it contributes no isolated candidate vertex. If it
has a solution \(x_I\), keep \(x_I\) exactly when it satisfies all inequalities
\(\langle a_i,x_I\rangle\leq1\). Duplicate candidates are identified. This
enumerates the vertex set of \(K\). At the same time it constructs the
vertex--facet incidence matrix
\[
  \mathcal I_{v i}
  =
  \begin{cases}
    1,&\langle a_i,v\rangle=1,\\
    0,&\langle a_i,v\rangle<1,
  \end{cases}
\]
where rows are vertices and columns are facets.

This incidence matrix is the concrete face-incidence data used later. A facet
is represented by the list of incident vertices in one column. Two facets
intersect if some row is incident to both. In four dimensions, an edge is
detected by two vertices sharing at least three incident facets, and a
two-face is detected by two facets sharing at least three vertices. These
are the incidence-only rules used by the current finite pipeline after the
vertex--facet matrix has been constructed exactly; they are not numerical
signed-distance tests. Degenerate inputs that need more than this incidence
contract are treated separately in the relevant algorithm section.

For generated dual rows, boundedness is checked before vertex enumeration.
Equivalently, the dual rows positively span \(\mathbb{R}^4\):
\[
  d\neq0
  \quad\Longrightarrow\quad
  \langle a_i,d\rangle>0
  \text{ for some }i.
\]
In four dimensions this can be checked from the one-dimensional kernels of
triples of dual rows, together with the rank condition that the \(a_i\) span
\(\mathbb{R}^4\). This is the boundedness test used by the exact polytope
pipeline.

We use the Hausdorff metric on compact convex sets for local and limiting
statements. Near a polytope with a fixed number of irredundant facets, a
labelled dual-row list gives local coordinates as long as the open conditions
above persist: \(0\in\operatorname{int}\operatorname{conv}\{a_i\}\), all
\(a_i\) remain extreme, and the relevant incidences do not change. Different
facet labellings can represent the same unlabelled polytope. This is the
topological background for local perturbation statements such as the
Hausdorff-local ten-facet chart in
Section~\ref{sec:hko-local-maximum}.

For a convex body \(K\), the support function and gauge function are
\[
  h_K(y) \coloneqq \max_{x\in K}\langle x,y\rangle,
  \qquad
  g_K(x) \coloneqq \inf\{\lambda>0:x\in \lambda K\}.
\]
The boundary of \(K\) is the level set \(g_K=1\). We use the Hamiltonian
\[
  H_K(x) \coloneqq g_K(x)^2.
\]
When \(K\) is smooth, a Reeb orbit on \(\partial K\) is the contact-normalized
Hamiltonian orbit
\begin{equation}\label{eq:preliminaries-smooth-reeb-inclusion}
  -J_0\dot{\gamma}(t) \in \partial H_K(\gamma(t)),
  \qquad
  H_K(\gamma(t))=1.
\end{equation}
Here \(\partial H_K\) denotes the convex subdifferential; in the smooth case it
is the singleton \(\{\nabla H_K\}\). With this normalization the action of a
Reeb orbit equals its period:
\[
  A(\gamma)=T.
\]

